\documentclass[12pt, babellanguage=british, titlespace]{altarticle}
\usepackage{bbmt}

\title{Documentation for the \texttt{bbmt} package}
\date{}

\newcommand{\optionline}[2][]{\par\medskip\hspace*{-.75cm}\texttt{#2}\hfill\if\relax\detokenize{#1}\relax\else\texttt{default:~#1}\fi\par}
\newcommand{\aargt}[1]{\ensuremath{\langle}\textrm{{\emph{#1}}}\ensuremath{\rangle}}
\newcommand{\argt}[1]{\texttt{\{}\ensuremath{\langle}\textrm{{\emph{#1}}}\ensuremath{\rangle}\texttt{\}}}
\newcommand{\optargt}[1]{\texttt{[}\ensuremath{\langle}\textrm{\emph{#1}}\ensuremath{\rangle}\texttt{]}}
\newcommand{\commandline}[1][]{\def\argI{#1}\ccommandline}
\newcommand{\ccommandline}[2][]{\par\medskip\hspace*{-.75cm}\texttt{\textbackslash#2}\argI{}\hfill\if\relax\detokenize{#1}\relax\else\texttt{default:~#1}\fi\par}
\newcommand{\pcommandline}[1][]{\def\argI{#1}\pccommandline}
\newcommand{\pccommandline}[2][]{\par\hspace*{-.75cm}\hphantom{\texttt{\textbackslash#2}}\argI{}\hfill\if\relax\detokenize{#1}\relax\else\texttt{default:~#1}\fi\par}
\newcommand{\envline}[1][]{\def\argI{#1}\eenvline}
\newcommand{\eenvline}[2][]{\par\medskip\hspace*{-.75cm}\texttt{\textbackslash begin\{#2\}}\argI{}~...~\texttt{\textbackslash end\{#2\}}\hfill\if\relax\detokenize{#1}\relax\else\texttt{default:~#1}\fi\par}

\begin{document}
\maketitle
This package provides a type 1 version of the \texttt{bbm} font. The glyphs were created using fontforge. It needs to be loaded after the \texttt{amsfonts} package.

\medskip

\commandline[\argt{characters}]{amsbb}This command outputs the blackboard font defined by the \ensuremath{\mathcal{A}\kern-.1667em\lower.5ex\hbox{\ensuremath{\mathcal{M}}}\kern-.125em\mathcal{S}}.

Here are the available characters (no bold version):

$\amsbb{ABCDEFGHIJKLMNOPQRSTUBWXYZ}$
\commandline[\argt{characters}]{mathbb}This command outputs the vecrotised version of the blackboard serif font defined by the \texttt{bbm} package.

Here are the available characters in normal weight:

$\mathbb{ABCDEFGHIJKLMNOPQRSTUBWXYZabcdefghijklmnopqrstuvwxyz12}$

Here are the available characters in bold weight:

{\boldmath$\mathbb{ABCDEFGHIJKLMNOPQRSTUBWXYZabcdefghijklmnopqrstuvwxyz12}$}
\commandline[\argt{characters}]{mathbbss}This command outputs the vecrotised version of the blackboard sans serif font defined by the \texttt{bbm} package.

Here are the available characters in normal weight:

$\mathbbss{ABCDEFGHIJKLMNOPQRSTUBWXYZabcdefghijklmnopqrstuvwxyz12}$

Here are the available characters in bold weight:

{\boldmath$\mathbbss{ABCDEFGHIJKLMNOPQRSTUBWXYZabcdefghijklmnopqrstuvwxyz12}$}
\end{document}